\section*{Guide to the Economic Applications}
This document is the current economic-applications companion to \emph{Neural Bellman Operators}. It retains the complete application text of the reviewed article and supplement, including model primitives, boundary and utility restrictions, analytical derivations, equilibrium conditions, implemented numerical procedures, and their recorded successes and failures. It is organized by economic problem. Its section, equation, table, and theorem numbers are distinct from those in the main article and technical supplement; cross-references connect all three documents.

The applications share a policy-evaluation and feasible-improvement construction. For a given policy, the critic solves that policy's continuation equation under the stated terminal, exit, or reflection conditions. The actor changes only the actions whose deviations define the relevant economic problem. A control problem compares feasible policies, a sophisticated consumer compares local deviations against future selves, and a firm compares complete unilateral deviations against fixed rivals. A small fitted residual is not substituted for any of these economic criteria.

\begin{table}[ht]
\centering
\caption{Economic objects in the application companion}
\label{tab:r16applicationsguide}
\begin{tabular}{p{0.22\textwidth}p{0.32\textwidth}p{0.35\textwidth}}
\toprule
Problem & Continuation object & Economic assessment\\
\midrule
Portfolio choice & Time-additive policy value & Analytical homothetic target and constrained actions\\
Recursive utility & The policy's own utility process & Utility domain, matching bequest, and comparison\\
Preference adjustment & Joint internal and wealth-state value & Consumption, portfolio, and adjustment choices; complete feasible-action accounts\\
Temporal selves & Two continuation components under a common future policy & Local spike-deviation equilibrium\\
Strategic investment & A distinct continuation value for each firm & Complete unilateral best response against fixed rivals\\
Linear-quadratic and capital laboratories & Coupled policy values and costates & Explicit representation tests, independent payoff and residual diagnostics\\
\bottomrule
\end{tabular}
\end{table}

The homothetic portfolio and recursive tests use a known invariant representation and report its scope explicitly. The endogenous-preference application retains the role of cardinal normalization across controlled preference states, the reflected preference domain, the cross derivative in portfolio choice, and the envelope comparison for adjustment costs. Its continuous-action records retain actor failures, local-search contributions, whole-domain off-grid failures, fresh-grid training, and the separate consumption-compensation experiment. The temporal-self condition concerns infinitesimal deviations; the game calculations retain missing pure fitted stage equilibria and approximate best-response bounds. The coupled linear-quadratic laboratory has an exactly representable structure. These distinctions are substantive economic restrictions rather than interchangeable numerical labels.

General approximation results and proofs are in the current technical supplement, beginning with Section~\ref{sec:theory}. The nonlinear capital economy and its current continuation-learning comparison are in the main article. The material below reproduces the complete application source with relocated input paths and documented grammar corrections; its reported quantities retain their original model, sample, and numerical meaning.
